\documentclass{article}

\usepackage{amsmath}
\usepackage{amsfonts}
\usepackage{amssymb}
\usepackage{xcolor}
\usepackage{listings}
\usepackage{graphicx}

\title{MS5610: Security Analysis and Portfolio Management}
\author{Lokesh Parihar}
\date{\today}

\begin{document}
\titlepage 
\newpage
\tableofcontents
\newpage

\section{28th July - Thursday}



\section{31st July - Friday}
\begin{itemize}
    \item 52 weeks high and low is a decent range
    \item what could be the value in the future?
    \item whether it is overvalued or undervalued against our future projections.
    \item 
\end{itemize}


\subsection{Security Valuation}
\begin{itemize}
    \item Debt Instruments
    \subitem Coupon Instruments
    \subitem Treasurey Instruments
    \subitem Zero Coupon Instruments
    \subitem Flexi Coupon Instruments
    \subitem Deep discount Instruments
    \item Equity Instruments
    \subitem Common Equity Instruments
    \subitem Preferred Equity Instruments
    \subitem Convertible Instruments
    \subitem Derivative Instruments

\end{itemize}

\subsection{Coupon Instruments}
\begin{itemize}
    \item LTP - Last Traded Price
    \item Interest on the coupons is not Traded
    \item Clean price, dirty price, and accrued Interest
    \item Trade happens at the clean price and accrued Interest is paid to the new buyer
    \item More commonly, the buyer pays the price deducing the accrued Interest
    \item 
\end{itemize}

\subsection{Valuation of Coupon Instruments}
\[\text{Running Yield} = \frac{Coupon amount}{Clean Price}\]

\[\text{Simple Yield to Maturity} = \frac{C}{CP} + \frac{P - CP}{n \cross CP}\]

\[\text{Price} = \sum_{i = 1}^{n} \frac{C_i}{(1 + YTM)^i} + \frac{P}{1 + YTM}^n\]

C - Coupon Amount, P - face value, CP - Clean Price, K - discount return

How do we find K? 

\[\text{Accrued Interest} = \frac{C}{m} \left(\frac{d_{xi} - d_{xc}}{T_d}\right)\]
$d_{xi}$ - number of days since the last coupon payment; 
$d_{xc}$ - number of days between book closure dates 

Term structure of Interest rates.

Exchange announces the date on which dividends are to be paid or something - Book Closure date (Please check it once)

ex- coupon; cum- coupon dates (Read about them)


\subsection{Equity Valuation: Multiplier Method}

\begin{itemize}
    \item Earnings Valuation
    \item Revenues Valuation
    \item Cash Flow Valuation
    \item Economic Value Added
    \item Accounting Value
    \item Yield Valuation
    \item Member's Valuation
\end{itemize}

\[\text{PE Multiplier} = \frac{\text{Market Price}}{\text{Earnings Per Share}}\]

\subsection{Multiplier Method}

Price-Eanings Ratio
\[PE Growth Rate = \frac{PE Multiplier}{Growth Rate}\]

\[PE Relative = \frac{Share PE}{Index PE}\]

Price to Sales Ratio
\[\frac{Market Capitalisation}{One year total revenue}\]

Price to Book Value Ratio 
\[\frac{Market Price}{Book value per share}\]



\section{04th August - Tuesday}
\subsection{Security Valuation}
\subsection{Class Exercise}
PSR = Price to Sales Ratio = Share Price/Revenue Per Share
PSR = Market Cap/Annual Revenue







\end{document}